\begin{figure}[ht]
\centering
\includegraphics[width=\textwidth]{figures/figS14_wiring_vs_compression.pdf}
\caption{\textbf{Supplementary Figure S13. Internalisation and compression are dissociable, not the same process.}
Each panel compares, on the same 12 model--prior $\times$ seed runs (plus the seed-matched
mismatch control R16), the prior-wiring effect size (per-edge cosine-coupling difference of the
injected edges, model vs the same seed's empty-prior model) against the latent-spectrum
compression ($\Delta\mathrm{PR}$, the participation ratio, $\Delta$ top-1\% concentration, $\Delta$ spectral
entropy). \textbf{(A--C)} The raw bivariate Spearman correlation between wiring and compression is
non-significant for $\Delta\mathrm{PR}$ ($\rho=-0.45$, $p=0.099$) and only nominally significant for
$\Delta$ entropy ($\rho=-0.54$, $p=0.034$); in both cases the apparent association is a
\emph{seed confound}: seed identity explains $72.7\%$ of the
variance in wiring and $83.0\%$ of the variance in $\Delta\mathrm{PR}$, and the
seed-controlled partial correlation (residualising each quantity by seed) collapses to zero
(Spearman $\rho=-0.07$ for $\Delta\mathrm{PR}$, $p=0.83$; $\rho=+0.01$ for $\Delta$ entropy,
$p=0.97$). The apparent
association is driven by seed, not by a mechanistic link. \textbf{(D)} The decisive dissociation:
the mismatched prior R16 compresses the marrow effective-rank spectrum almost as much as the true
prior ($\Delta\mathrm{PR}=-2.89$ versus R13's $-3.44$, $\sim$84\%) while wiring
\emph{neither} edge set ($\Delta$ cosine $-0.034$, $z=-1.43$, n.s.). Compression therefore occurs
without targeted internalisation. \textbf{Conclusion:} latent-spectrum compression is a
bulk-spectral consequence of injecting \emph{any} structured GRN prior, whereas internalisation of
the prior's specific TF$\rightarrow$target edges is a separate, program-matching-gated event that
requires the prior to match the trajectory's active drivers. The two should not be used
interchangeably when interpreting prior-injected VAEs.}
\label{fig:wiring_vs_compression}
\end{figure}
